\documentclass[11pt,a4paper]{article}
\usepackage[utf8]{inputenc}
\usepackage{graphicx}
\usepackage[left=2.5cm,top=3cm,right=2.5cm,bottom=3cm,bindingoffset=0.5cm]{geometry}
\usepackage{AEDLogica, AEDEspecificacion, AEDTADs}
\usepackage{caratula}

\newcommand{\Tipo}[1]{\mathsf{#1}}
% la sintaxis es \newcommand{\nombreDeLaMacro}[cantidadDeParametros]{Lo que va ser remplazado por el macro} 
\newcommand{\norm}[1]{\vert#1\vert}



\begin{document}

\section{Renombres de Tipos y Estructuras}

\obs{Usuarios}{\dict{id:\Tipo{\Z}}{Categoria}}

\noindent \comentario{La convencion utilizada para transacciones es:
    \\
    1) Si se hace un vuelo, el TipoDeTransaccion es 0 y el usuario de destino es -1
    \\
    2) Si se hace una transferencia el TipoDeTransaccion es 1
    \\
    3) Si se hace un canjeo el TipoDeTransaccion es 2 y el usuario de destino es -1, se modifica unicamente el usuario de origen
}
\\
\noindent Transacciones \textbf{ES} Struct $\tupla{
        UsuarioOrigen: \Tipo{\Z}, \\
        UsuarioDestino: \Tipo{\Z},
        TipoDeTransaccion: \Tipo{\Z},
        Millas: \Tipo{\R}, \\
        Distancia: \Tipo{\Z},
        NombrePromo: \Tipo{\str},
        FactorPromo: \Tipo{\R}
    }$

\noindent \obs{Historial}{\seq{Transacciones}}
\\
\comentario{En caso de las promociones, si la promocion no tiene un tope, por convencion tomamos millasRestantes = -1}
\\
\texttt{obs} Promociones : dict \{ \\
\hspace*{1.5em} NombrePromocion: $\str$, \\
\hspace*{1.5em} $\struct{\Tipo{FactorDePromocion} : \R,\Tipo{topeDePromocion}: \R, \Tipo{millasRestantes}: \R}$ \\
\}


\begin{tad}{Airmiles}

    % ===============Aca comienza la funcion de acumular millas===============

    \noindent \comentarioLargo{Acumula millas al usuario si la categoria del mismo no es platino, en caso de que el nombre de la promo no alcance para completar el viaje del usuario, se crean dos transacciones, una con las millas restantes de la promo y otra con "NoPromo" para completar el viaje. }
    \begin{proc}{AcumularMillas}
        {
            \Inout AM: \Tipo{Airmiles},
            \In idUsuario: \Tipo{\Z},
            \In distancia: \Tipo{\R},
            \In nombrePromo: \Tipo{\str},
        }
        {}

        \requiere{AM = AM0 \land idUsuario \in claves(AM.Usuarios) \land nombrePromo \in claves(AM.Promociones) \land distancia > 0 \land
            AM.Usuarios[idUsuario].categoria \ne Platino}
        \aseguraLargo{
            \comentario{Si el tope de las millas de la promo supera la distancia se hace una sola transaccion}
            \\
            AM0.Promociones\lbrack nombrePromo \rbrack.millasRestantes \geq
            \\
            MillasCalculadas(distancia, DarFactorCategoria(AM.Historial,idUsuario),
            \\ AM.Promociones[nombrePromo].factorDePromocion) \implicaLuego
            \\
            \\
            EsNuevaTransaccion(AM,AM0,idUsuario,-1,0,
            \\
            MillasCalculadas(distancia, DarFactorCategoria(AM.Historial,idUsuario),
            \\ AM.Promociones[nombrePromo].factorDePromocion) ,distancia, \\
            nombrePromo, AM0.Promociones[nombrePromo].FactorDePromocion) \lor
            \\
            \\
            AM0.Promociones\lbrack nombrePromo \rbrack.millasRestantes < distancia \implica
            \\
            \\
            \comentario{Se hace 1 transaccion con las millas restantes de promo}
            \\
            EsNuevaTransaccion(AM,AM0,idUsuario,-1,1,
            \\
            AM0.Promociones[nombrePromo].millasRestantes, distancia,
            \\ nombrePromo, AM0.Promociones[nombrePromo].FactorDePromocion) \yLuego
            \\
            \\
            \comentario{Se hace otra transaccion con noPromo para completar}
            \\
            EsNuevaTransaccion(AM,AM0,idUsuario,-1,1,
            \\
            millasACompletar(AM0.Promociones,nombrePromo,distancia,
            \\DarFactorCategoria(AM0.Historial,idUsuario),
            \\AM0.Promociones
                [nombrePromo].FactorDePromocion), \\distancia, "NoPromo", 1) \yLuego
            % Luego de hacer la transaccion tengo que verificar si se modifico la categoría del usuario
            \\
            \comentario{Se puede indefinir si no le pasamos un usuario que este en el sistema}
            \\
            SubioCategoria(idUsuario,AM.Historial,AM0.Historial) \implicaLuego
            \\
            AM.Usuarios = SetKey(AM0.Usuarios,idUsuario,
            \\
            DarCategoria(AM.Historial, idUsuario)) \land
            \\
            \comentario{Despues de completar el viaje se verifica si la promocion termino}
            \\
            SeConsumioLaPromocion(nombrePromo,AM0.Promociones)
        }
    \end{proc}

    \auxLargo{MillasObtenidasPorVuelo}
    {Historial: \Tipo{\seq{Transacciones}}, idOrigen: \Tipo{\Z}}
    {\Tipo{\R}}
    {
        \sum\limits_{i=0}^{\norm{Historial} - 1}
        \IfThenElse{
            EsVuelo(Historial[i],
            idOrigen)
        }{
            Historial[i].Millas
        }{
            0
        }
    }

    \vspace{1cm}

    \auxLargo{MillasTransferidas}
    {Historial: \Tipo{\seq{Transacciones}}, idOrigen: \Tipo{\Z}, idDestino: \Tipo{\Z}}
    {\Tipo{\R}}
    {
        \sum\limits_{i=0}^{\norm{Historial} - 1}
        \IfThenElse{
            EsTransferencia(Historial[i],
            \\
            idOrigen, idDestino)
        }{
            Historial[i].Millas
        }{
            0
        }
    }

    \vspace{1cm}
    \auxLargo{MillasRecibidas}
    {AM: \Tipo{Airmiles}, idDestino: \Tipo{\Z}}
    {\Tipo{\R}}
    {
        \sum\limits_{i=0}^{\norm{AM.Historial} - 1}
        \IfThenElse{
            EsTransferenciaRecibida(AM.Historial[i],
            \\
            idDestino, AM.Usuarios)
        }{
            AM.Historial[i].Millas
        }{
            0
        }
    }

    \vspace{1cm}
    \predLargo{EsTransferenciaRecibida}
    {Transaccion: \Tipo{Transaccion}, idDestino: \Tipo{\Z}, Usuarios: \Tipo{\dict{\Z}{Categoria}}}
    {
        Transaccion.UsuarioOrigen \in claves(Usuarios) \land
        \\
        Transaccion.UsuarioOrigen \neq idDestino
        \\
        Transaccion.TipoDeTransaccion = 1 \land
        \\
        Transaccion.UsuarioDestino = idDestino \land
        \\
        \comentario{Deberia verificar los otros campos? Se complica...}
    }

    \vspace{1cm}
    \noindent \comentario{En caso de que el usuario haga parte del viaje con promo y otra parte sin promo esta funcion busca esa diferencia}
    \\
    \auxLargo{MillasACompletar}
    {Promociones: \Tipo{\Diccionario{\str}{struct}},
        \\nombrePromo: \Tipo{\str}, distancia: \Tipo{\Z}, FactorCategoria: \Tipo{\R}, FactorPromo: \Tipo{\R}}
    {\R}
    {
        MillasCalculadas(distancia,factorCategoria,factorPromo)-
        \\
        Promociones[nombrePromo].millasRestantes
    }

    \vspace{1cm}
    \auxLargo{MillasDescontadasPorCanjeo}
    {Historial: \Tipo{\seq{Transacciones}}, idOrigen: \Tipo{\Z}}
    {\Tipo{\R}}
    {
        \sum\limits_{i=0}^{\norm{Historial} - 1}
        \IfThenElse{
            EsCanjeo(Historial[i],
            idOrigen)
        }{
            Historial[i].Millas
        }{
            0
        }
    }

    \vspace{1cm}

    \predLargo{EsVuelo}
    {Transaccion: \Tipo{Struct}, idOrigen: \Tipo{\Z}}
    {
        Transaccion.UsuarioOrigen = idOrigen \land
        \\
        Transaccion.TipoDeTransaccion = 0 \land
        \\
        Transaccion.UsuarioDestino = -1 \land
        \\
        \comentario{Deberia verificar los otros campos? Se complica...}
    }

    \vspace{1cm}

    \predLargo{EsTransferencia}
    {Transaccion: \Tipo{struct}, idOrigen: \Tipo{\Z}, idDestino: \Tipo{\Z}}
    {
        Transaccion.UsuarioOrigen = idOrigen \land
        \\
        Transaccion.TipoDeTransaccion = 1 \land
        \\
        Transaccion.UsuarioDestino = idDestino \land
        \\
        Transaccion.millas > 0
        \\
        \comentario{Deberia verificar los otros campos? Se complica...}
    }

    \vspace{1cm}



    \predLargo{EsCanjeo}
    {Transaccion: \Tipo{struct}, idOrigen: \Tipo{\Z}}
    {
        Transaccion.UsuarioOrigen = idOrigen \land
        \\
        Transaccion.TipoDeTransaccion = 2 \land
        \\
        Transaccion.UsuarioDestino = -1 \land
        \\
        Transaccion.distancia > 0
        \\
        \comentario{Deberia verificar los otros campos? Se complica...}
    }

    \vspace{1cm}
    \noindent \comentario{Recordar que aca tambien estan incluidos los referidos porque esta incluido en las "MillasObtenidasPorVuelo"}
    \\
    \auxLargo{MillasTotales}
    {Historial: \Tipo{\seq{Transacciones}}, idOrigen: \Tipo{\Z}, idDestino: \Tipo{\Z}}
    {\Tipo{\R}}
    {
        MillasObtenidasPorVuelo(Historial, idOrigen) +
        \\
        \comentario{en este orden son las que me transfirieron}
        \\
        MillasTransferidas(Historial, idDestino, idOrigen)
    }

    \vspace{1cm}
    \auxLargo{MillasDisponibles}
    {AM: \Tipo{Airmiles}, idOrigen: \Tipo{\Z}, idDestino: \Tipo{\Z}}
    {\Tipo{\R}}
    {
        MillasTotales(AM.Historial, idOrigen) -
        \\
        MillasDescontadasPorCanjeo(AM.Historial,idOrigen) -
        \\
        \comentario{en este orden son las que transferi (descontadas)}
        \\
        MillasTransferidas(AM.Historial, idOrigen, idDestino) +
        \\
        MillasRecibidas(AM.Historial,idOrigen, AM.Usuarios)
    }

    \vspace{1cm}
    \auxLargo{DarCategoria}
    {Historial: \Tipo{\seq{Transacciones}}, idUsuario: \Tipo{\Z}}
    {\Tipo{Categoria}}
    {
        \IfThenElse{
            MillasTotales(Historial, idUsuario) < 10000
        }{
            Bronce
            \\
        }{
            \IfThenElse{
                MillasTotales(Historial, idUsuario) < 50000
            }{
                Plata
                \\
            }{
                \IfThenElse{
                    MillasTotales(Historial, idUsuario) < 1000000
                }{
                    Oro
                }{
                    Platino
                }
            }
        }
    }

    \vspace{1cm}
    \auxLargo{DarFactorCategoria}
    {Historial: \Tipo{\seq{Transacciones}}, idUsuario: \Tipo{\Z}}
    {\Tipo{\R}}
    {
        \IfThenElse{
            MillasTotales(Historial, idUsuario) < 10000
        }{
            1.0
            \\
        }{
            \IfThenElse{
                MillasTotales(Historial, idUsuario) < 50000
            }{
                1.5
                \\
            }{
                \IfThenElse{
                    MillasTotales(Historial, idUsuario) < 1000000
                }{
                    2.0
                }{
                    0.0
                }
            }
        }
    }

    \vspace{1cm}
    \auxLargo
    {MillasCalculadas}
    {distancia:\Tipo{\Z}, Factorcategoria: \Tipo{\R}, factorPromo: \Tipo{\R}}
    {\R}
    {distancia * Factorcategoria * factorPromo}

    \vspace{1cm}
    \noindent \comentario{Chequea si el usuario cambio de categoria con la lista de transacciones que se crearon hasta ese momento y con la nueva lista de transacciones}
    \\
    \predLargo
    {SubioCategoria}
    {idUsuario: \Tipo{\Z},Historial: \Tipo{Transaccion},Historial0: \Tipo{Transaccion}}
    {DarCategoria(Historial0,idUsuario) \neq DarCategoria(Historial,idUsuario)}

    \predLargo
    {SeConsumioLaPromocion}
    {
        nombrePromo: \Tipo{\str},
        Promociones: \Tipo{\Diccionario{\str}{struct}},
        \\
        Promociones0: \Tipo{\Diccionario{\str}{struct}}}
    {AM.Promociones = DelKey(AM.Promociones0,nombrePromo)}

\end{tad}

\end{document}
% Hay que chequear varias cosas:
% proc AcumularMillas(inout AM: Airmiles, in: idUsuario:Z, in: distancia:R, in: nombrePromo:str){
% Requiere{ AM = AM0 ∧ idUsuario ∈   claves(AM.Usuarios) ∧ nombrePromo ∈ claves(AM.Promociones)}
% Asegura{

% ActualizaMillasUsuario(idUsuario, Usuarios, Usuarios0,distancia,FactorCategoria,FactorPromo) ∧L

% SubioCategoria(idUsuario,Usuarios) ∧ 

% EsNuevaTransaccion(AM, AM0, idUsuario, idReferidor, TipoDeTransaccion, CalcularMillas(distancia,FactorCategoria,FactorPromo,idUsuario, Usuarios),
% NombrePromo, FactorPromo)} ∧L

% SeConsumioLaPromocion(nombrePromo,Promociones)

% }